\scanpage{81}\mspage{77}
Ainsi, les spatios peuvent être considérés comme des disjoncteurs particuliers, qu'on pourrait appeler \uline{disjoncteurs} \uline{canoniques}. \deleted{\unsure{Si}}

Si $(E,\inserted{\mathrel{|\unsure{\circ|0}|}})$ est un disjoncteur quelconque, on aimerait considérer $ \mathbb{\Sigma}_{E}$, muni de
\[
\sigma_{E}:E\unsure{-|\longrightarrow}\mathbb{\Sigma}_{E},
\]
comme un ``disjoncteur canonique'' (ou \unsure{Spatios|spatios}) \uline{enveloppant} de $E$. \unsure{Je voudrais|Je voudraits} déjà caractériser, à is. unique près, le spatios $\mathbb{\Sigma}_{E}$ et l'application $\sigma$.

\uline{Proposition 2} \unsure{Soient|Soit} $(E,\mathrel{|\circ|})$ un disjoncteur, $(\deleted{\unsure{S|\mathscr{S}}}\Sigma,\leq,C)$ un spatios, et
\[
\alpha:E\longrightarrow\Sigma
\]
une application. Pour qu'il existe un diagramme commutatif
\diagram{\begin{tikzcd}[ampersand replacement=\&,row sep=0.6em,column sep=3em]
{} \& \deleted{\unsure{E|S}}\Sigma \\
\unsure{,}\,E \arrow[ur,"{\deleted{\unsure{\alpha|\varphi}}\sigma}"] \arrow[dr,"{\sigma_E}"] \& {} \\
{} \& \mathbb{\Sigma}_E \arrow[uu,leftrightarrow,"{\rotatebox{90}{$\sim$}\ \substack{\deleted{\unsure{\alpha|\varphi}}\\ \unsure{\varphi|1}}}"{swap}] \& {,}
\end{tikzcd}}
où $\varphi$ est un iso, il f. et s. que $\sigma$ satisfasse les \inserted{deux} conditions suivantes
\[
\forall x,y\in E,\qquad x\mathrel{|\circ|}y\Longleftrightarrow
\deleted{\unsure{\alpha|\varphi}}\sigma x\mathrel{|\circ|}\deleted{\unsure{\alpha|\varphi}}\sigma y
\tag*{\makebox[0pt][r]{\makebox[\displaywidth][l]{a)}}}
\]
% archive sheet number: 80
